\documentclass[../../../include/open-logic-section]{subfiles}

\begin{document}

\olfileid[tr]{sth}{spine}{zf}
\olsection{$\Z$ ve $\ZF$: Bir Dönüm Noktası}

Temellendirme ile bir başka önemli dönüm noktasına ulaşırız. Belirli bir şey
``çıkarılmış'' olarak tanımladığımız $\Zminus$ ve $\ZFminus$ teorilerini
ele aldık. Bu belirli şey Temellendirme'dir. Öyleyse:

\begin{defn}
$\Z$ teorisi $\Zminus$'a Temellendirme'yi ekler. Dolayısıyla aksiyomları
Kaplamsallık, Birleşim, Çiftler, Kuvvet Kümeleri, Sonsuzluk, Temellendirme ve
Ayırma şemasının bütün örnekleridir.

$\ZF$ teorisi $\ZFminus$'a Temellendirme'yi ekler. Başka bir deyişle
$\ZF$, $\Z$'ye Yerine Koyma'nın bütün örneklerini ekler.
\end{defn}

Yine de aklınıza bir soru gelmiş olabilir. Düzenlilik $\ZFminus$ üzerinde
Temellendirme'ye denkse ve Düzenlilik'in gerekçesi açıksa, neden dolambaçlı
bir yol izleyip temel aksiyomumuz olarak Düzenlilik yerine Temellendirme'yi
alıyoruz?

Tarihsel nedenleri (kimin neyi ne zaman formüle ettiğiyle ilgili nedenleri)
bir yana bırakırsak temel neden, Temellendirme'nin $V_\alpha$'ların tanımını
kullanmadan sunulabilmesidir. Bu tanım \olref[recursion]{sec} içindeki
bütün çalışmaya dayanıyordu: gerekçeli olduğunu göstermek için Sonlu Ötesi
Özyineleme'yi ispatlamamız gerekti. Fakat Sonlu Ötesi Özyineleme ispatımız
\emph{Yerine Koyma}'yı kullandı. Dolayısıyla Temellendirme ile Düzenlilik
$\ZFminus$'a göre denk olsa da $\Zminus$'a göre denk değildir.

Aslında durum bu basit açıklamanın düşündürdüğünden daha çarpıcıdır. Bu
kitabın kapsamını epey aşsa da hem $\Zminus$ hem de $\Z$'nin
$V_\alpha$'ları tanımlayamayacak kadar zayıf olduğu ortaya çıkar. Bu yüzden
yalnızca $\Z$ içinde çalışıyorsanız Düzenlilik (bizim formüle ettiğimiz
biçimiyle) bir \emph{anlam} bile taşımaz. Resmî aksiyomumuzun Düzenlilik
değil Temellendirme olmasının nedeni budur.

Bundan sonra, aksi belirtilmedikçe ve başka bir açıklama yapmadan, $\ZF$
içinde çalışacağız.

\end{document}
